\section{Results}
\label{sec:results}

\paragraph{Overview of the integrated knowledge graph}
\label{sec:knowledge_graph}

We constructed an integrated knowledge graph to analyse the representation
and interoperability of regulatory substances across datasets.
Figure~\ref{fig:knowledge_graph} provides an overview of the resulting
graph structure and its main semantic layers.

\begin{figure}[!ht]
    \centering
    \includemermaid[width=1\linewidth]{images/sources/knowledge_graph.mmd}
    \caption{
        Overview of the integrated knowledge graph constructed from 18 regulatory datasets.
        Regulatory substances (~17,500 entities) are represented as SKOS concepts
        and form the core layer of the graph.
        Where possible, substances are linked to structure-based identifiers (InChIKey),
        enabling alignment with chemical classification systems (ChemOnt)
        and biological roles (ChEBI).
        In addition, substances are organised into linkability tiers
        reflecting their degree of chemical specificity, and into
        semantic clusters derived from embedding-based similarity.
        The graph is further enriched with annotations (~46,600 links)
        connecting substances to regulatory sources and functional roles.
        The figure highlights the multiple semantic layers
        required to represent and integrate heterogeneous regulatory substance definitions.
    }\label{fig:knowledge_graph}
\end{figure}

\paragraph{Dataset Overview}

Integration of 18 regulatory sources yielded 41,813 records. Of these,
63\% could be linked to a defined molecular structure via InChIKey.

%Table~\ref{tab:datasources} provides an overview of all regulatory and monitoring sources included in this study.

%\input{../output/tables/datasources}

\paragraph{Heterogeneity of Substance Definitions}
\label{sec:semantic_mismatch}

A significant share of regulatory entries (37\%) does not correspond to a defined chemical structure, consistent with Kutsarova et al.~\cite{Kutsarova2019}. These include mixtures (UVCBs), substance groups, analytical parameters, or entries without usable identifiers, reflecting differences in how “substances” are defined across regulatory contexts.

Heterogeneity varies widely between regulatory lists: some are mostly structure-based, while others rely heavily on group or mixture definitions, limiting interoperability. (Figure~\ref{fig:regulatory_groups})

Applying the seven-tier linkability taxonomy (Section~\ref{sec:linkability}) shows that 62.5\% of entries (20,909) are directly linkable via InChIKey (tier 1). Only a small fraction consists of partially resolvable groups (tiers 2–3; ${<}0.2$\%).

The remaining 37.5\% cannot be directly linked to a structure: 11.6\% have unresolved CAS numbers (tier 4), 10.8\% are mixtures or UVCBs (tier 5), ${<}0.1$\% are administrative entries (tier 6), and 14.9\% lack any usable identifier (tier 7).

Overall, more than one-third of regulatory entries cannot be mapped to a unique chemical structure, and even CAS numbers do not guarantee unambiguous identification. Consequently, structure-based integration applies only to part of the data, requiring alternative approaches for the remainder.

\begin{figure}[H]
    \centering
    \includegraphics[width=\linewidth]{Analysis_4bd_v_Entity_type_and_linkability_vertical.pdf}
    \caption{
        (A) Entity type composition across 18 regulatory sources. Each bar shows the proportional breakdown of substance entries into seven mutually exclusive types (Molecule, Substance group, Analytical parameter, Mixture, Regulatory entry, CAS without structure, Unclassified), illustrating that the mix of entity types differs substantially between regulatory instruments. (B) Linkability taxonomy applied to all 33,452 unique regulatory substance entries. A total of 20,909 entries (62.5\%) can be directly linked to a chemical structure via InChIKey (tier 1), while 4,981 entries (14.9\%) lack any usable identifier. Together, the panels show that structural linkability is unevenly distributed across sources and that more than one-third of all entries cannot be resolved to a unique chemical structure.
    }
    \label{fig:analysis4d}
\end{figure}

\paragraph{Limitations of CAS Numbers}

Although CAS numbers are widely used, they do not provide a reliable
unique identifier. Multiple CAS numbers may refer to the same structure,
and conversely, a single CAS number may correspond to multiple structures.
This confirms that CAS-based linking is insufficient for interoperability.

\paragraph{Limited Interoperability Across Regulatory Lists}

Cross-list analysis shows that most substances are unique to a single
regulatory source. Overlap between lists is limited, reflecting the
domain-specific design of regulatory instruments rather than a shared
chemical reference framework.

\paragraph{Failure of Semantic Approaches for Non-Structure Substances}
Embedding-based clustering of non-structure substance names yields semantically coherent groups, but does not resolve chemical identity. While clusters reflect lexical or contextual similarity between substance descriptions, they do not correspond to well-defined molecular entities (Figure~\ref{fig:analysis89}A).

\begin{figure}[H]
    \centering
    \includegraphics[width=\linewidth]{Analysis_89_UMAP_and_ChemOnt_quality.pdf}
    \caption{(A) UMAP projection of sentence embeddings for 12,499 non-structure regulatory substance entries, coloured by k-means cluster (k = 6). Each point represents a unique substance name; spatial proximity reflects semantic similarity in the 384-dimensional embedding space. (B) Distribution of top-1 cosine similarity scores~\cite{Manning2008} between substance name embeddings and ChemOnt class label embeddings across all 12,499 candidate substance-class pairs, shown here for \texttt{all-MiniLM-L6-v2}. The dashed line indicates the applied threshold ($s_1 > 0.80$); 133 entries (1.1\%) are retained as matches.}
    \label{fig:analysis89}
\end{figure}

To evaluate whether semantic similarity can recover chemical meaning, we attempted to map non-structure substances to ChemOnt taxonomy classes using cosine similarity between sentence embeddings of substance names and ChemOnt class labels. Four models were evaluated: \texttt{all-MiniLM-L6-v2} (miniLM), \texttt{all-mpnet-base-v2} (mpnet), \texttt{allenai/scibert\_scivocab\_uncased} (SciBERT), and \texttt{dmis-lab/biobert-v1.1} (BioBERT). The distribution of top-1 similarity scores is heavily skewed towards lower values (Figure~\ref{fig:analysis89}B). Cosine similarity scores are not comparable across models: score floors at the 200th-ranked match range from approximately 0.80 for miniLM to 0.94 for BioBERT, reflecting differences in embedding geometry rather than differences in match quality. A fixed threshold such as $s_1 > 0.80$ is therefore model-dependent and arbitrary; it was applied here for the purpose of bounding a manual evaluation, not as a principled quality criterion.

Manual inspection of the top-200 matches per model revealed that the majority of retained assignments are chemically incorrect, regardless of model. Correct assignments were largely limited to cases where the substance name closely resembles the ChemOnt class label (e.g.\ \textit{Phenols} $\to$ \textit{Phenols}, \textit{Alkyl mercury compounds} $\to$ \textit{Alkyl mercury compounds}), indicating that high similarity scores reflect surface-form overlap rather than semantic understanding. Systematic errors were present across all models: miniLM conflates organic and inorganic species and confuses elements sharing a common prefix (e.g.\ \textit{organostannic compounds} $\to$ \textit{organoarsenic compounds}); mpnet consistently misassigns all nickel sulfide variants to \textit{nitrile sulfides}, a token-level artefact of the chemical symbol Ni being read as the start of ``nitrile''; SciBERT and BioBERT, pre-trained on biomedical literature, exhibit severe score inflation---a score of 0.93 on an incorrect match provides less information than a score of 0.85 on a correct one---and produce repeated false anchors to specific heterocyclic and lipid ChemOnt classes for structurally unrelated substances. Trade names and reaction products were assigned arbitrary classes with no verifiable basis across all models.

Critically, score magnitude is not a reliable indicator of match correctness. Models with the highest score floors (SciBERT, BioBERT) performed worst on both manual inspection and quantitative validation, demonstrating that inflated cosine scores reflect poor calibration rather than genuine chemical understanding. Embedding-based methods cannot compensate for the absence of structure-based identifiers, and fail systematically for substance groups, mixtures, and analytically defined entries.

\paragraph{Validation of Embedding-Based ChemOnt Classification}
\label{sec:chemont_results}

To quantify the ability of text-based embeddings to recover chemical
taxonomy assignments from substance names, embedding models were evaluated
against a structure-derived ChemOnt ground truth. The validation set
comprised 28,204 structure-defined substances with ClassyFire-based
assignments to the ChemOnt hierarchy.

Model performance was assessed using standard information retrieval metrics,
including Hit@k and Mean Reciprocal Rank (MRR). Among the evaluated models,
\texttt{all-MiniLM-L6-v2} consistently outperformed the others, achieving a
Hit@1 of 5.75\%, Hit@3 of 9.86\%, Hit@5 of 12.34\%, and Hit@10 of 16.21\% for
the direct parent ChemOnt class. The corresponding MRR was 0.093, indicating
that the correct class was typically ranked around the top ten positions.
In contrast, \texttt{all-mpnet-base-v2}, \texttt{allenai/scibert\_scivocab\_uncased},
and \texttt{dmis-lab/biobert-v1.1} showed substantially lower performance,
with Hit@1 values below 1\% and MRR values below 0.02.

A hierarchy-aware evaluation further revealed that retrieval performance
decreased with increasing taxonomic abstraction. For the best-performing
model, Hit@1 was 3.25\% at the \emph{subclass} level and below 1\% for
\emph{class}, \emph{superclass}, and \emph{kingdom}, indicating limited
semantic alignment between substance names and higher-level ChemOnt
categories. A weighted lenient MRR, designed to account for partial matches
across the hierarchy, yielded a value of 0.099 for the MiniLM model, only
marginally higher than the strict MRR.

Detailed performance metrics for all evaluated models are provided in
Table~\ref{tab:chemont_validation}.

\begin{table}[!ht]
    \centering
    \caption{Performance of embedding models for ChemOnt classification.}
    \label{tab:chemont_validation}
    \begin{tabular}{lccccc}
        \hline
        Model           & Hit@1  & Hit@3  & Hit@5  & Hit@10 & MRR   \\
        \hline
        all-MiniLM-L6-v2  & 0.0575 & 0.0986 & 0.1234 & 0.1621 & 0.093 \\
        all-mpnet-base-v2 & 0.0093 & 0.0152 & 0.0184 & 0.0247 & 0.015 \\
        SciBERT           & 0.0040 & 0.0071 & 0.0102 & 0.0131 & 0.008 \\
        BioBERT           & 0.0047 & 0.0077 & 0.0099 & 0.0150 & 0.009 \\
        \hline
    \end{tabular}
\end{table}

\paragraph{Scale of the Interoperability Problem}
For the more than 15,000 non-structure entries, full manual reconciliation would require on the order of hundreds of millions of pairwise comparisons. Under conservative assumptions, this corresponds to several thousand person-years of expert effort.

\label{sec:workload}
\begin{figure}[!ht]
    \centering
    \includegraphics[width=\linewidth]{Analysis_10c_workload_non_structured.pdf}
    \caption{
        Estimated pairwise curation workload for 11,250 non-structure entries across the combined European regulatory dataset, broken down by source.
        The total corresponds to approximately 630,000 person-days (3,150 person-years) under an assumed throughput of 100 pair assessments per person-day.
        The quadratic growth of the workload illustrates that manual reconciliation is not scalable, reinforcing the need for structure-based and ontology-driven representations at the source.
    }
    \label{fig:analysis10c}
\end{figure}



Figure~\ref{fig:analysis10c} quantifies this workload across all regulatory sources, illustrating both the absolute scale and the quadratic growth that makes each additional non-structure entry disproportionately costly. Even moderate increases in the number of unresolved entries lead to a rapidly escalating curation burden.

Manual pairwise reconciliation at this scale is not operationally feasible. This result highlights a fundamental constraint: interoperability cannot be achieved through retrospective harmonisation of heterogeneous substance definitions.

Instead, the only scalable strategy is to address the problem at its source. By requiring structure-resolvable identifiers and ontology-based representations in regulatory definitions wherever chemically applicable, the number of entries requiring manual reconciliation can be structurally reduced rather than managed post hoc.

\paragraph{Knowledge Graph Integration}

Structure-defined substances can be successfully integrated across
regulatory lists, chemical taxonomies (ChemOnt), and biological
annotations (ChEBI) using InChIKey as a linking key. This demonstrates
the feasibility of a linked-data approach when structure information is
available.

\paragraph{Priority Patterns}

The integrated knowledge graph enables the combination of regulatory
and biological annotations into a shared analytical space.
Using the composite priority score defined in Section~\ref{sec:priority_score},
we identify substances and chemical classes that are prominent across
both regulatory presence and biological role annotations.

The bubble chart (Figure~\ref{fig:analysis14e}) visualises this
priority space by plotting $\text{list\_score}$ against
$\text{bio\_score}$ for all scored substances.
Bubble size is proportional to the number of regulatory source memberships.


\begin{figure}[!ht]
    \centering
    \includegraphics[width=\linewidth]{Analysis_14e_Bubble_priority.pdf}
    \caption{Bubble chart of composite priority space: regulatory list score versus biological hazard score for all scored substances. Bubble size is proportional to the number of regulatory source memberships. Substances in the upper-right quadrant combine broad regulatory coverage with documented biological hazard evidence.}
    \label{fig:analysis14e}
\end{figure}

Substances located in the upper-right quadrant combine high regulatory
coverage with multiple biological role annotations, reflecting their
prominence across heterogeneous data sources.
Examples include bisphenol~A (BPA) and dibutyl phthalate, which appear
consistently across regulatory lists and are associated with multiple
biological roles in ChEBI.

Importantly, this analysis does not constitute a toxicological ranking,
but demonstrates how interoperable data representations enable the
integration of regulatory and ontology-based information within a
single analytical framework.

At the class level, a small number of chemical classes accounts for
a large fraction of regulated substances. This concentration suggests
that class-based representations can provide a scalable bridge between
individual substances and regulatory group definitions, supporting
alignment across regulatory contexts.




